%!TEX root = ./main.tex

\section{ソースコード}

この節では，listings を使ったソースコードの表示を例示します．

\subsection{C++}

\begin{pre}
\begin{lstlisting}[language=C++]
#include <iostream.h>

void main() {
    cout << "Hello World!" << endl;
    return 0;
}
\end{lstlisting}
\end{pre}

\subsection{Haskell}

\begin{pre}
\begin{lstlisting}[language=Haskell]
hamming = 1 : map (*2) hamming # map (*3) hamming # map (*5) hamming
    where xxs@(x:xs) # yys@(y:ys)
              | x==y = x : xs # ys
              | x<y  = x : xs # yys
              | x>y  = y : xxs # ys
\end{lstlisting}
\end{pre}

\subsection{HTML}

\begin{pre}
\begin{lstlisting}[language=html]
<!DOCTYPE html>
<html>
  <head>
    <title>This is a title</title>
  </head>
  <body>
    <p>Hello world!</p>
  </body>
</html>
\end{lstlisting}
\end{pre}

\subsection{Java}

\begin{pre}
\begin{lstlisting}[language=Java]
package fibsandlies;
import java.util.HashMap;

/**
 * This is an example of a Javadoc comment; Javadoc can compile documentation
 * from this text. Javadoc comments must immediately precede the class, method, or field being documented.
 */
public class FibCalculator extends Fibonacci implements Calculator {
    private static Map<Integer, Integer> memoized = new HashMap<Integer, Integer>();

    /*
     * The main method written as follows is used by the JVM as a starting point for the program.
     */
    public static void main(String[] args) {
        memoized.put(1, 1);
        memoized.put(2, 1);
        System.out.println(fibonacci(12)); //Get the 12th Fibonacci number and print to console
    }

    /**
     * An example of a method written in Java, wrapped in a class.
     * Given a non-negative number FIBINDEX, returns
     * the Nth Fibonacci number, where N equals FIBINDEX.
     * @param fibIndex The index of the Fibonacci number
     * @return The Fibonacci number
     */
    public static int fibonacci(int fibIndex) {
        if (memoized.containsKey(fibIndex)) {
            return memoized.get(fibIndex);
        } else {
            int answer = fibonacci(fibIndex - 1) + fibonacci(fibIndex - 2);
            memoized.put(fibIndex, answer);
            return answer;
        }
    }
}
\end{lstlisting}
\end{pre}

\subsection{Python}

\begin{pre}
\begin{lstlisting}[language=Python]
import random
class Die(object):
#A dice has a feature of number about how many sides it has when it's established,like 6.
def __init__(self):
  self.sides=6

"""because a dice contains at least 4 planes.
So use this method to give it a judgement when you need to change the instance attributes."""
def set_sides(self, sides_change):
  if sides_change>=4:
    if sides_change != 6:
      print("change sides from 6 to ",sides_change," !")
    else:  # added else clause for printing a message that sides set to 6
        print ("sides set to 6")
    self.sides = sides_change
  else:
    print("wrong sides! sides set to 6")

def roll(self):
  return random.randint(1, self.sides)

d = Die()
d1 = Die()
d.set_sides(4)
d1.set_sides(4)
print (d.roll(), d1.roll())
\end{lstlisting}
\end{pre}
